\chapter{Penutup}

\section{Kesimpulan}

Penelitian ini menyelidiki struktur teori gravitasi pada ruang-waktu tiga dimensi (dua dimensi ruang dan satu dimensi waktu) melalui tiga tahap: pembuktian ekuivalensi antara gravitasi Einstein dan teori gauge Chern-Simons, penurunan solusi lubang hitam BTZ, dan analisis sektor Einstein-Maxwell. Temuan-temuan utama dari masing-masing tahap dirangkum berikut ini.

Ekuivalensi antara gravitasi Einstein-Palatini dan teori Chern-Simons berhasil dibuktikan secara eksplisit melalui ansatz Witten, yang menggabungkan koneksi spin dan dreibein menjadi satu koneksi gauge tunggal bernilai pada aljabar Lie isometri ruang anti-de Sitter. Seluruh struktur aljabar, termasuk relasi komutasi antar-generator rotasi Lorentz dan translasi AdS, diturunkan langsung dari prinsip pertama tanpa mengasumsikan bentuk aljabar secara \textit{ad hoc}. Persamaan medan yang dihasilkan dari variasi aksi Chern-Simons berupa syarat koneksi datar; setelah didekomposisi ke dalam basis generator aljabar, syarat ini mereproduksi secara tepat dua persamaan fundamental gravitasi: kondisi bebas torsi yang menentukan koneksi spin sebagai fungsi dreibein, dan persamaan medan Einstein vakum dengan konstanta kosmologi negatif. Hasil ini mengonfirmasi bahwa gravitasi tiga dimensi, yang tidak memiliki derajat kebebasan propagasi lokal, memang sepenuhnya bersifat topologis dan dapat diformulasikan sebagai teori gauge.

Solusi lubang hitam BTZ diturunkan secara analitik dari persamaan medan Einstein menggunakan formalisme dreibein. Dengan ansatz ruang-waktu yang stasioner dan aksisimetris, penyelesaian kondisi torsi-nol menentukan seluruh komponen koneksi spin sebagai fungsi dari parameter metrik. Syarat kelengkungan konstan (yang merupakan realisasi persamaan Einstein vakum pada tiga dimensi) kemudian menentukan fungsi-fungsi metrik secara unik dalam dua parameter: massa dan momentum sudut. Metrik yang dihasilkan mendeskripsikan lubang hitam dengan dua horizon (horizon peristiwa luar dan dalam) yang posisinya ditentukan oleh kedua parameter tersebut. Kasus statik diperoleh sebagai limit momentum sudut nol, di mana horizon dalam menghilang dan metrik mereduksi ke lubang hitam BTZ tanpa rotasi.

Perluasan analisis ke sektor Einstein-Maxwell, di mana medan elektromagnetik terkopling ke gravitasi melalui tensor energi-momentum Maxwell, menghasilkan tiga temuan. Pertama, persamaan Maxwell pada latar belakang melengkung dapat diselesaikan secara umum tanpa mengasumsikan bentuk fungsi metrik terlebih dahulu; komponen medan listrik dan medan magnetik dinyatakan dalam parameter muatan dan fungsi metrik yang belum ditentukan. Kedua, terdapat kendala geometris fundamental yang melarang keberadaan lubang hitam BTZ yang sekaligus bermuatan listrik dan berotasi dalam teori Einstein-Maxwell murni tiga dimensi. Akar fisis dari kendala ini terletak pada ketidakcocokan antara simetri geometri dan simetri sumber: ansatz dreibein BTZ memaksa kelengkungan pada dua komponen tertentu menjadi identik (suatu kesamaan struktural yang berlaku untuk sembarang fungsi metrik), sementara tensor energi-momentum Maxwell merespons secara asimetris pada kedua komponen tersebut. Satu-satunya cara agar persamaan Einstein bersumber terpenuhi secara konsisten adalah melenyapkan sumber ketidakcocokan tersebut, yakni muatan listrik nol (kembali ke BTZ vakum) atau momentum sudut nol (lubang hitam statik). Kendala ini bukan artefak pemilihan gauge atau penyederhanaan teknis, melainkan konsekuensi intrinsik dari struktur geometri ansatz BTZ yang digabungkan dengan sifat tensor energi-momentum Maxwell. Ketiga, satu-satunya solusi bermuatan yang konsisten adalah lubang hitam statik, yang fungsi metriknya mengandung suku logaritmik sebagai manifestasi dari potensial Coulomb dua dimensi ruang. Horizon peristiwanya ditentukan oleh persamaan transendental yang mencampur suku polinomial dan logaritmik, sehingga tidak memiliki solusi analitik tertutup, berbeda dari kasus BTZ vakum yang horizonnya diperoleh dari persamaan aljabar.


\section{Saran}

Beberapa arah pengembangan yang dapat ditempuh berdasarkan hasil penelitian ini:

\begin{enumerate}[label=\textbf{\arabic*.}, leftmargin=*, itemsep=6pt]
    \item \textbf{Kopling medan dilaton.} Kendala eksistensi lubang hitam bermuatan dan berotasi yang ditemukan dalam penelitian ini bersifat spesifik untuk teori Einstein-Maxwell murni, di mana medan elektromagnetik terkopling secara minimal ke gravitasi. Penambahan medan skalar dilaton $\phi$ dengan kopling eksponensial $f(\phi) = e^{-2\alpha\phi}$ ke suku kinetik Maxwell dapat mengubah struktur tensor energi-momentum sedemikian rupa sehingga asimetri yang menjadi sumber kontradiksi tidak lagi muncul. Chan dan Mann \cite{chan1994} memperoleh keluarga solusi lubang hitam bermuatan statik pada teori Einstein-Maxwell-dilaton $2+1$ dimensi, di mana aksinya terkait secara konformal dengan aksi efektif string energi-rendah. Baru-baru ini, Karakasis \textit{et al.}\ \cite{karakasis2023} berhasil memperoleh solusi lubang hitam berotasi eksak pada teori yang sama; parameter dilaton $\alpha$ mempengaruhi fungsi metrik, sifat rotasi, dan termodinamika lubang hitam, serta mereduksi ke BTZ biasa pada limit $\alpha \to 0$. Kedua referensi tersebut menggunakan formalisme metrik (notasi tensor). Mengerjakan ulang teori Einstein-Maxwell-dilaton ini dalam formalisme bentuk diferensial dan vielbein, sebagaimana digunakan dalam penelitian ini, merupakan pengembangan yang natural: suku kopling dilaton-Maxwell diterjemahkan menjadi $e^{-2\alpha\phi}\, F \wedge {*}F$, suku kinetik dilaton menjadi $d\phi \wedge {*}d\phi$, dan potensial dilaton masuk sebagai $V(\phi)\,\eta$. Reformulasi ini diharapkan tidak hanya mereproduksi solusi berotasi-bermuatan yang telah diketahui, tetapi juga membuka kemungkinan untuk memformulasikan sektor dilaton dalam kerangka Chern-Simons yang diperluas.

    \item \textbf{Korespondensi AdS/CFT dan entropi mikroskopis.} Formulasi Chern-Simons yang dibangun dalam penelitian ini menyediakan kerangka kerja yang natural untuk menghitung muatan permukaan pada batas asimtotik ruang anti-de Sitter tiga dimensi melalui prosedur Brown-Henneaux. Analisis kondisi batas asimtotik dapat memunculkan aljabar konformal Virasoro dengan muatan sentral tertentu, yang kemudian memungkinkan perhitungan entropi mikroskopis lubang hitam BTZ melalui formula Cardy. Verifikasi kesesuaian antara entropi mikroskopis ini dan entropi termodinamika Bekenstein-Hawking akan memberikan uji konsistensi non-trivial bagi ekuivalensi yang telah dibuktikan.
\end{enumerate}